\section{Reproducibility, Successes, and Challenges}

This section documents how the project is rebuilt, where the historical
replication succeeds, and where source definitions or data revisions create
material differences. The report contains no code snippets; implementation
and executable examples belong in the repository and notebook tour.

\subsection{End-to-End Rebuild}

The project separates raw-data acquisition, tidy-data construction, empirical
analysis, and reporting. PyDoit orchestrates these stages and generates all
numerical table entries, figure files, captions, metadata, and the final
LaTeX report. The final build records the data vintage, sample endpoints, Git
commit, and artifact checksums.

\subsection{Replication Audit}

\input{../tables/table_r1_replication_audit.tex}

\subsection{Successes and Remaining Challenges}

The machine-generated audit summary is \ReplicationAuditSummary. Strict passes
identify close historical matches; revised-vintage passes isolate results that
remain consistent after allowing the predeclared revision tolerance; failures
remain visible for diagnosis. The primary risk-free convention is
\RiskFreePrimary{}, and the primary term-spread convention is
\TermSpreadPrimary{}. The pipeline constructs and compares both candidates
before recording those choices. Changes in updated coefficients are reported
as new empirical evidence, not as historical replication failures.

\subsection{Data and Copyright Boundary}

The repository records citations, series identifiers, transformations, and
checksums, but does not commit copyrighted papers, screenshots, licensed CRSP
extracts, or secrets. Local reference materials are resolved outside Git.
